\chapter{トークナイザーの作成}
\label{ch:tokenizer}

この章では、テキストを数値シーケンスに変換する\keyterm{トークナイザー（tokenizer)}を直接実装します。 2つの主なアルゴリズムであるBPE（Byte Pair Encoding）とUnigramを扱い、それぞれの学習、エンコーディング、デコーディングの過程をコードで見ていきます。

\section{なぜトルクナイザーが必要なのか}

ニューラルネットワークは数字だけを理解する。 「Hello、world！」という文字列をそのままモデルに入れることはできません。テキストを整数シーケンスに変換する最初のステップが必要です。これを\keyterm{トークナイザー}が担当します。最も簡単な方法は「文字単位のトークナイザー」で、すべての文字を別々のトークンとして扱うことです。しかし、これはシーケンスが長すぎて単語の意味をつかむのが難しいです。

反対の極端は「単語単位トークナイザー」であり、スペースで区切られた単語を1つのトークンとして扱います。しかし、これは語彙サイズが爆発し（英語だけでも数十万語）、未学習語（out-of-vocabulary、OOV）を処理できません。

現代LLMは、これら2つの極端の中間である\keyterm{サブワード（subword)}トルクナイザーを使用しています。よく登場する単語はそのまま1つのトークンで、珍しい単語はより小さなサブワードに分割する。 「unhappiness」は、「un」+「happiness」または「un」+「happi」+「ness」に分割することができます。このアプローチは、語彙サイズを適切に維持しながらOOV問題を解決します。

\subsection{3つのアプローチの比較}

\begin{center}
\begin{tabularx}{\textwidth}{lXXX}
\toprule
\textbf{方法}&\textbf{語彙サイズ}&\textbf{OOV 処理}&\textbf{シーケンス長}\\
\midrule
レタリングユニット＆小（$\sim$100）＆なし＆超キム\\
単語単位＆非常に大きい（$\sim$100K）＆処理不能＆短い\\
サブワード＆中間（$\sim$10K$\sim$50K）＆部分単語に分解＆中間\\
\bottomrule
\end{tabularx}
\end{center}

\section{特殊トークン}

すべてのトルクナイザーはいくつかの\keyterm{特殊トークン（special token)}を持っています。私たちの実装では4つの方法を使用します。

\begin{center}
\begin{tabularx}{\textwidth}{llX}
\toprule
\textbf{トークン}&\textbf{ID}&\textbf{用途}\\
\midrule
\code{<pad>}＆0＆バッチ内のシーケンスの長さを合わせるためのパディング\\
\code{<bos>}&1 & シーケンスの開始を表示 (Beginning Of Sequence) \\
\code{<eos>}& 2 & シーケンスの終わりを表示 (End Of Sequence) \\
\code{<unk>}＆3＆語彙にないトークン（Unknown）\\
\bottomrule
\end{tabularx}
\end{center}

これらのトークンはインデックスが固定されており、モデルが特別に扱うことができる。たとえば、生成時に\code{<eos>}が出た場合、推論を停止し、\code{<pad>}位置はloss計算で無視します。

\begin{warnbox}[特殊トークンID固定の理由]
特殊トークンのIDを0$\sim$3に固定すると、モデルコードが単純化されます。 「\code{if token\_id == 0: mask\_out}」のようなハードコーディングが可能です。また、チェックポイントの互換性が保証されています。トークナイザーを再学習しても、特別なトークンIDは常に同じであるため、モデルの重みをそのまま再利用できます。
\end{warnbox}

\section{ベースクラス}

すべてのトークナイザーが従うべきインタフェースを抽象クラスとして定義します。

\begin{lstlisting}[language=Python]
from abc import ABC, abstractmethod
from dataclasses import dataclass

@dataclass
class SpecialTokens:
    pad: str = "<pad>"   # ID 0
    bos: str = "<bos>"   # ID 1
    eos: str = "<eos>"   # ID 2
    unk: str = "<unk>"   # ID 3

    @property
    def tokens(self) -> list[str]:
        return [self.pad, self.bos, self.eos, self.unk]

    @property
    def pad_id(self) -> int: return 0
    @property
    def bos_id(self) -> int: return 1
    @property
    def eos_id(self) -> int: return 2
    @property
    def unk_id(self) -> int: return 3


class Tokenizer(ABC):
    def __init__(self, special_tokens: SpecialTokens | None = None):
        self.special = special_tokens or SpecialTokens()
        self.id_to_token: list[str] = list(self.special.tokens)
        self.token_to_id: dict[str, int] = {t: i for i, t in enumerate(self.id_to_token)}

    @abstractmethod
    def _train(self, corpus: list[str], vocab_size: int) -> None:
        """語彙学習（子クラスで実装）"""

    @abstractmethod
    def _tokenize(self, text: str) -> list[str]:
        """テキストをトークンリストに分割する（子クラスで実装）"""

    def encode(self, text: str, add_bos=False, add_eos=False) -> list[int]:
        tokens = self._tokenize(text)
        ids = [self.token_to_id.get(t, self.special.unk_id) for t in tokens]
        if add_bos: ids = [self.special.bos_id] + ids
        if add_eos: ids = ids + [self.special.eos_id]
        return ids

    def decode(self, ids: list[int]) -> str:
        tokens = [self.id_to_token[i] for i in ids
                  if self.id_to_token[i] not in self.special.tokens]
        raw = self._tokens_to_bytes(tokens)
        return self._postprocess(raw.decode("utf-8", errors="replace"))
\end{lstlisting}

\begin{infobox}[抽象クラスの役割]
\code{Tokenizer}は「インタフェース」を定義します。子クラス（BPE、Unigram）は\code{\_train}と\code{\_tokenize}を実装する必要があります。これにより、学習アルゴリズムは異なりますが、\code{encode/decode}APIは同じままで、上位コード（データセット、トレーナー）がトルクナイザーの種類に関係なく動作します。

これが「多形性」の価値です。トレーナーは\code{tokenizer.encode(text)}を呼び出すだけで、内部的にBPEが回っても Unigramが回っても気にしない。後で「SentencePiece」や「WordPiece」などの新しいアルゴリズムを追加しても、トレーナーコードを変更する必要はありません。
\end{infobox}

\section{BPE トークナイザー}

\keyterm{BPE(Byte Pair Encoding)}は、1994年にデータ圧縮のために提案されたアルゴリズムを、2016年にSennrichらがニューラルネットワークの機械翻訳に導入したものです。 GPT-2、GPT-3、LLaMAなどが使用する。

\subsection{アルゴリズム直観}

BPEのコアアイデアは単純です。
\begin{enumerate}
\item テキストをバイト単位で分解する（すべての文字を表現可能、OOVなし）。
\item 隣接するバイトペアの中で最も頻繁に登場するペアを探します。
\item そのペアを新しいトークンにマージします。
\item 目標語彙サイズに達するまで2$\sim$3を繰り返します。
\end{enumerate}

最初はすべてのテキストが「バイトのリスト」です。たとえば、「ローワー」は\code{[l, o, w, e, r]}で表されます。学習を進めながら：
\begin{itemize}
\item 「lo」が頻繁に登場すると、マージ：\code{[lo, w, e, r]}
\item 「er」が頻繁に登場すると、マージ：\code{[lo, w, er]}
\item 「low」が頻繁に登場すると、マージ：\code{[low, er]}
\item 「Lower」が頻繁に登場すると、マージ：\code{[lower]}
\end{itemize}

最後に、「ローワー」は1つのトークンになりますが、まれな単語は部分単語として残ります。 「Lowering」は「Lower」+「ing」になります。

\begin{figure}[htbp]
\centering
\begin{tikzpicture}[
  every node/.style={font=\small\sffamily},
  byte/.style={draw=inkblue, fill=inkblue!8, minimum width=0.8cm, minimum height=0.6cm, rounded corners=2pt},
  merge/.style={draw=inkred, fill=inkred!10, minimum width=1.6cm, minimum height=0.6cm, rounded corners=2pt},
  arr/.style={-{Stealth[length=4pt]}, inkgray}
]
  \node[byte] (a1) at (0, 0) {l};
  \node[byte, right=0.1cm of a1] (a2) {o};
  \node[byte, right=0.1cm of a2] (a3) {w};
  \node[byte, right=0.1cm of a3] (a4) {e};
  \node[byte, right=0.1cm of a4] (a5) {r};
  \node[inkgray, anchor=west, font=\scriptsize] at (5.5, 0) {step 0: bytes};

  \node[merge] (b1) at (0, -1) {lo};
  \node[byte, right=0.1cm of b1] (b2) {w};
  \node[byte, right=0.1cm of b2] (b3) {e};
  \node[byte, right=0.1cm of b3] (b4) {r};
  \node[inkgray, anchor=west, font=\scriptsize] at (5.5, -1) {step 1: merge 'l'+'o'};

  \node[merge] (c1) at (0, -2) {low};
  \node[byte, right=0.1cm of c1] (c2) {e};
  \node[byte, right=0.1cm of c2] (c3) {r};
  \node[inkgray, anchor=west, font=\scriptsize] at (5.5, -2) {step 2: merge 'lo'+'w'};

  \node[merge] (d1) at (0, -3) {low};
  \node[merge, right=0.1cm of d1] (d2) {er};
  \node[inkgray, anchor=west, font=\scriptsize] at (5.5, -3) {step 3: merge 'e'+'r'};

  \node[merge, fill=inkblue!20, draw=inkblue, line width=0.8pt] (e1) at (0, -4) {lower};
  \node[inkgray, anchor=west, font=\scriptsize] at (5.5, -4) {step 4: merge 'low'+'er'};

  \foreach \y in {0, -1, -2, -3} {
    \draw[arr] (\y+0.05, \y+0.3) -- (\y+0.05, \y+0.7);
  }
\end{tikzpicture}
\caption{BPE 学習コース. ``lower''がバイトで始まり、徐々にマージされて1つのトークンになるプロセス.}
\label{fig:bpe-process}
\end{figure}

\subsection{なぜバイト単位なのか?}

GPT-2以降、BPEは「バイト単位」で動作します。以前は文字単位を書きましたが、問題がありました。
\begin{itemize}
\item 絵文字、珍しい漢字など語彙にない文字は\code{<unk>}処理されます。
\item 多言語コーパスで文字セットが大きすぎます。
\end{itemize}

バイト単位の場合、すべてのUTF-8文字を1  --  0 4バイトで表現でき、語彙の「バイト」部分は常に256に固定されています。 OOVは事実上消えます。

\subsection{実装: 学習}

\begin{lstlisting}[language=Python]
from collections import Counter
from .base import Tokenizer

class BPETokenizer(Tokenizer):
    def __init__(self, special_tokens=None):
        super().__init__(special_tokens=special_tokens)
        self.merges: list[tuple[str, str]] = []
        self._merge_rank: dict[tuple[str, str], int] = {}

    def _train(self, corpus: list[str], vocab_size: int) -> None:
        # 1. 特殊トークン + 256バイトで初期語彙を構成する
        self.id_to_token = list(self.special.tokens)
        for b in range(256):
            self.id_to_token.append(chr(b))
        self.token_to_id = {t: i for i, t in enumerate(self.id_to_token)}

        # 2. コーパスを単語単位に分割
        word_freqs = Counter()
        for text in corpus:
            for word in self._pre_tokenize(text):
                word_freqs[word] += 1

        # 3. 各単語をバイトトークンリストに変換
        word_to_tokens = {
            w: [chr(b) for b in w.encode("utf-8")]
            for w in word_freqs
        }

        # 4. マージループ
        self.merges = []
        target = vocab_size - len(self.id_to_token)
        for _ in range(max(target, 0)):
            pair_freqs = Counter()
            for word, freq in word_freqs.items():
                toks = word_to_tokens[word]
                for i in range(len(toks) - 1):
                    pair_freqs[(toks[i], toks[i+1])] += freq
            if not pair_freqs:
                break
            best = max(pair_freqs.items(), key=lambda x: (x[1], x[0]))
            if best[1] < 2:
                break
            new_token = best[0][0] + best[0][1]
            self.merges.append(best[0])
            self.id_to_token.append(new_token)
            self.token_to_id[new_token] = len(self.id_to_token) - 1
            for word in word_to_tokens:
                word_to_tokens[word] = self._merge_in_list(
                    word_to_tokens[word], best[0], new_token
                )

        self._merge_rank = {p: i for i, p in enumerate(self.merges)}
\end{lstlisting}

\begin{notebox}[マージ優先順位の安定性]
\code{max(pair\_freqs.items(), key=lambda x: (x[1], x[0]))}から\code{(x[1], x[0])}タプルにソートする理由は、同点処理によるものです。同じ頻度のペアが複数ある場合、アルファベット順に最前線のペアを選択して結果を決定論的にします。こうしてこそ同じコーパスで学習すると、常に同じマージシーケンスが出る。
\end{notebox}

\subsection{エンコーディング}

学習したマージルールを新しいテキストに適用するには、最も優先順位の高い（最初に学習した）マージから適用します。

\begin{lstlisting}[language=Python]
    def _tokenize(self, text: str) -> list[str]:
        tokens = []
        for word in self._pre_tokenize(text):
            bs = word.encode("utf-8")
            word_tokens = [chr(b) for b in bs]
            word_tokens = self._apply_merges(word_tokens)
            tokens.extend(word_tokens)
        return tokens

    def _apply_merges(self, tokens: list[str]) -> list[str]:
        """学習したマージ規則を優先順位に適用する."""
        while len(tokens) >= 2:
            # 最も rank低い（優先順位が高い）ペアを探す
            best_idx, best_rank = -1, len(self.merges) + 1
            for i in range(len(tokens) - 1):
                rank = self._merge_rank.get((tokens[i], tokens[i+1]), -1)
                if 0 <= rank < best_rank:
                    best_rank, best_idx = rank, i
            if best_idx < 0:
                break
            merged = tokens[best_idx] + tokens[best_idx + 1]
            tokens = tokens[:best_idx] + [merged] + tokens[best_idx + 2:]
        return tokens
\end{lstlisting}

\subsection{前処理: スペース処理}

BPEはスペースを特別に扱わなければ単語の境界を保存できない。 GPT-2はスペースを特殊文字Ġ（U + 0120）に置き換えます。私たちはSentencePieceスキームに従ってU + 2581（）に置き換えます。

\begin{lstlisting}[language=Python]
    def _pre_tokenize(self, text: str) -> list[str]:
        text = text.replace(" ", "▁")
        text = text.replace("\n", "▁").replace("\t", "▁")
        while "▁▁" in text:
            text = text.replace("▁▁", "▁")
        if not text.startswith("▁"):
            text = "▁" + text
        parts = text.split("▁")
        return ["▁" + p for i, p in enumerate(parts) if i > 0 and p]
\end{lstlisting}

\begin{warnbox}[空白置換の罠]
「Hello world」を「Hello world」に置き換えてからスプリットすると\code{["▁Hello", "▁world"]}になります。単語ごとにこの前に付いており、デコード時を再び空白に変えると元の文章が復元される。

しかし、「Hello、world！」のように句読点がある場合は注意が必要です。単純なスプリットは「Hello」を作成し、カンマが単語の一部になります。実際のGPT-2は正規表現パターンで単語/句読点を分離しますが、この本では単純化のために省略します。
\end{warnbox}

\subsection{デコード: バイトの復元}

BPEのトリッキーな部分はデコードです。トークンは\code{chr(byte)}のシーケンスなので、各文字のコードポイントをバイトとして扱い、元のバイトストリームを復元してからUTF-8にデコードする必要があります。

\begin{lstlisting}[language=Python]
    def _tokens_to_bytes(self, tokens: list[str]) -> bytes:
        """BPE トークンをバイトに変換. 各文字のコードポイントをバイト単位で."""
        out = bytearray()
        for tok in tokens:
            for c in tok:
                out.append(ord(c) & 0xFF)
        return bytes(out)
\end{lstlisting}

たとえば、「He​​llo」トークンがある場合、（U + 2581）はUTF-8で\code{0xE2 0x96 0x81}なので3バイトですが、BPE学習時に「」文字自体を3つの別々のバイトトークンとして扱ったので\code{chr(0xE2) + chr(0x96) + chr(0x81)}で表されます。デコード時に各\code{chr(b)}から\code{ord()}にバイトを抽出すると元のバイトストリーム\code{0xE2 0x96 0x81 0x48 0x65 0x6C 0x6C 0x6F}が得られ、それをUTF-8にデコードすると「Hello」になります。

\section{Unigram トークナイザー}

\keyterm{Unigram}トルクナイザーは、Kudo（2018）がSentencePieceの一部として提案した方法です。 BPEが「頻度ベースのマージ」の場合、Unigramは「確率ベースの分割」です。

\subsection{直観}

Unigramは各トークンに確率を割り当てます。文章「ローワー」について考えられる分割はいくつかあります。
\begin{itemize}
\item\code{[l, o, w, e, r]}：5つのトークン
\item\code{[lo, w, er]}：3つのトークン
\item\code{[low, er]}：2つのトークン
\item\code{[lower]}：1トークン
\end{itemize}

各分割の確率はトークン確率の積として計算される。最も高い確率を持つ分割を選択します。 「低」が語彙にあり、確率が高い場合は\code{[lower]}が選択されますが、存在しない場合は小さい単位に分割されます。

\subsection{アルゴリズム}

\begin{enumerate}
\item コーパスからすべての部分文字列を候補トークンに抽出し、頻度を数えます。
\item 頻度を確率に変換して、各トークンにログ確率スコアを割り当てます。
\item EM（Expectation-Maximization）反復：
    \begin{itemize}
\item E-step：Viterbiアルゴリズムによる各単語の最適分割ナビゲーション。
\item M-step: 分割結果でトークン確率を更新。
    \end{itemize}
\item 最もスコアの低いトークンを削除し、語彙を縮小します。
\end{enumerate}

\subsection{Viterbi デコード}

与えられた単語に対して「最も高い総ログ確率を持つトークン分割」を動的計画法で見つけます。

\begin{lstlisting}[language=Python]
    def _viterbi(self, word: str) -> tuple[list[str], float]:
        n = len(word)
        # dp[i] = (最大スコア, 前のインデックス, 以前のトークン）
        dp = [(-float("inf"), -1, "") for _ in range(n + 1)]
        dp[0] = (0.0, -1, "")
        for i in range(1, n + 1):
            for j in range(max(0, i - 16), i):
                substr = word[j:i]
                score = self.token_scores.get(substr, -20.0)
                cand = dp[j][0] + score
                if cand > dp[i][0]:
                    dp[i] = (cand, j, substr)
        # 逆追跡
        path, i = [], n
        while i > 0:
            _, j, tok = dp[i]
            path.append(tok)
            i = j
        path.reverse()
        return path, dp[n][0]
\end{lstlisting}

\begin{infobox}[Viterbiの複雑さ]
$O(n^2)$の代わりに$O(n \cdot L)$に縮小された秘密は\code{max(0, i-16)}部分です。トークンの長さが16を超えることはまれであるため、$j$の範囲を制限して検索スペースを減らします。実際のSentencePieceも同様の最適化を使用します。
\end{infobox}

\section{BPE vs Unigram 比較}

\begin{center}
\begin{tabularx}{\textwidth}{lXX}
\toprule
 & \textbf{BPE} & \textbf{Unigram} \\
\midrule
学習方法＆頻度ベースのマージ＆確率ベースの分割（EM）\\
エンコード＆マージルールの順序を適用＆Viterbiで最適な分割\\
決定論的＆O（同じ入力、同じ出力）＆O \\
マルチスプリットサポート&X&O(サンプリング可能) \\
主な用途 & GPT-2, GPT-3, LLaMA & T5, ALBERT, mBART \\
\bottomrule
\end{tabularx}
\end{center}

\section{トルクナイザーテスト}

実装が終わったら、単体テストで検証します。最も重要なテストは、\keyterm{ラウンドトリップ（round-trip) テスト}：エンコード後にデコードすると原文を復元する必要があるということです。

\begin{lstlisting}[language=Python]
def test_encode_decode_roundtrip():
    tok = BPETokenizer()
    tok.train(CORPUS, vocab_size=200)
    text = "the quick fox"
    ids = tok.encode(text)
    decoded = tok.decode(ids)
    assert " ".join(decoded.split()) == "the quick fox"

def test_bos_eos_addition():
    tok = BPETokenizer()
    tok.train(CORPUS, vocab_size=100)
    ids = tok.encode("hello", add_bos=True, add_eos=True)
    assert ids[0] == tok.special.bos_id
    assert ids[-1] == tok.special.eos_id

def test_unknown_token_handled():
    """語彙にない文字もバイト単位で処理されます。 unkが発生しないこと."""
    tok = BPETokenizer()
    tok.train(CORPUS, vocab_size=100)
    ids = tok.encode("こんにちは")
    assert tok.special.unk_id not in ids  # バイトに分解されるので unk なし

def test_save_load(tmp_path):
    tok = BPETokenizer()
    tok.train(CORPUS, vocab_size=150)
    path = tmp_path / "tok.json"
    tok.save(path)
    tok2 = BPETokenizer.load(path)
    text = "the quick fox"
    assert tok.encode(text) == tok2.encode(text)
\end{lstlisting}

\begin{notebox}[テスト哲学]
この本のすべてのモジュールには3つのレベルのテストがあります。
\begin{itemize}
\item\textbf{形状テスト（shape test)}：テンソルの形状が予想と一致するかどうか。
\item\textbf{数値テスト（numeric test)}: 小さい入力に対する手計算結果と比較。
\item\textbf{統合テスト}: 複数のモジュールを接続したときのエンドツーエンドの動作。
\end{itemize}
これら3つを通過したコードだけが本に収録される。

ラウンドトリップテストは統合テストの一例です。 「エンコード後にデコードすると原文が出る」ということは自明に見えますが、実際には空白処理、マルチバイト文字、特殊トークンの削除など、さまざまな段階でバグが発生する可能性があります。このテストに合格すると、パイプライン全体が正しく動作することを確認できます。
\end{notebox}

\section{練習: 実際の馬でトルクナイザーを学ぶ}

本の\code{scripts/tiny\_train.py}デモを実行すると、実際にBPEトークナイザーを学習するプロセスを見ることができます。

\begin{lstlisting}[style=bashstyle]
cd make-llm-basic
python scripts/tiny_train.py --vocab 500
\end{lstlisting}

出力例：
\begin{lstlisting}[basicstyle=\ttfamily\footnotesize, backgroundcolor=\color{codebg}, frame=single]
[1/5] トークナイザー学習中...
    完了: 308トークン, 1.2秒
    サンプル: 'the quick brown fox jumps over the lazy dog'
    → 18トークン: [264, 282, 291, 267, 261, 110, ...]
\end{lstlisting}

\section{要約}

\begin{itemize}
\item トルクナイザーは、テキストを整数シーケンスに変換するLLMの最初のステップです。
\item サブワード方式は語彙サイズとOOVのバランスをとります。
\item BPEは頻度ベースのマージ、Unigramは確率ベースの分割。
\item 特殊トークン\code{<pad>/<bos>/<eos>/<unk>}はID 0$\sim$3に固定されています。
\item バイト単位のBPEはOOVを事実上除去します。
\item ViterbiアルゴリズムでUnigramの最適分割を$O(n)$に探します。
\item すべての実装はpytestユニットテストで検証されています。
\end{itemize}

\section{練習問題}

\begin{enumerate}
\item BPEを学習するとき、「最も頻度の高いペア」の代わりに「最も頻度の低いペア」をマージするとどうなりますか？結果を予測して実験してください。
\item Unigramトルクナイザーで\code{score = -20.0}ペナルティ値を\code{-5.0}に変更するとどうなりますか？
\item 韓国語コーパス「こんにちはお会いできて嬉しいです」をBPEとして学びましょう。語彙サイズが100のときと1000のときにどのような違いがありますか？
\item Viterbiでトークン長の上限を16ではなく4に減らすと、どのトークンが消えますか？
\item\code{<bos>}トークンがなぜ必要なのか考えてください。なくてもモデルが動作するか？
\end{enumerate}

次の章では、トークンIDを密集ベクトルに変換する埋め込みレイヤを実装します。
